\section{Обучение моделей классификации на данных, полученных в предыдущих НИРС}

\subsection{Подготовка данных к обучению}

Исходным материалом для данного этапа служит датасет, сформированный в рамках НИРС-02~\cite{github_repo_nirs2}. Датасет содержит 8\,115 Bitcoin-адресов, каждый из которых описывается 34 признаками. Признаки сгруппированы по смыслу:
\begin{itemize}
    \item \textit{структурные}~--- число входов и выходов транзакций, количество уникальных контрагентов, средняя вентиляторность (fan-in/fan-out);
    \item \textit{UTXO-паттерны}~--- доля нераспределённых выходов, средний остаток на адресе, характеристики типичного изменения UTXO;
    \item \textit{временны́е}~--- общий период активности адреса, средний интервал между транзакциями, число активных дней;
    \item \textit{экономические}~--- суммарный объём входящих и исходящих средств (в satoshi), средний объём одной транзакции, соотношение объёмов;
    \item \textit{Bitcoin-специфичные}~--- наличие SegWit-выходов, участие в мультиподписных схемах, паттерны округления сумм.
\end{itemize}

Целевая переменная~--- тип кошелька~--- принимает шесть значений: майнеры~(\textit{miner},~3\,547 адресов, 43,7\%), биржи~(\textit{exchange},~1\,532, 18,9\%), сервисы~(\textit{services},~1\,249, 15,4\%), азартные игры~(\textit{gambling},~911, 11,2\%), операции типа CoinJoin~(\textit{coinjoin-like},~787, 9,7\%) и майнинговые пулы~(\textit{mining\_pool},~89, 1,1\%).

\subsubsection{Отбор признаков и устранение мультиколлинеарности}

Наличие высококоррелированных признаков в обучающей выборке нежелательно: оно не несёт дополнительной информации, но увеличивает размерность пространства, замедляет обучение и может нестабилизировать оценки важности признаков. Степень линейной зависимости между парой непрерывных признаков $x_i$, $x_j$ измеряется коэффициентом корреляции Пирсона:
\[
    r(x_i, x_j) = \frac{\sum_{k=1}^{n}(x_i^{(k)} - \bar{x}_i)(x_j^{(k)} - \bar{x}_j)}
    {\sqrt{\sum_{k=1}^{n}(x_i^{(k)} - \bar{x}_i)^2 \cdot \sum_{k=1}^{n}(x_j^{(k)} - \bar{x}_j)^2}},
\]
где $x_i^{(k)}$, $x_j^{(k)}$~--- значения признаков $x_i$ и $x_j$ для $k$-го образца; $\bar{x}_i$, $\bar{x}_j$~--- их выборочные средние; $n$~--- число образцов. Если $|r(x_i, x_j)| > 0{,}95$, один из признаков удаляется. По результатам автоматического анализа корреляционной матрицы исключены четыре признака: \textit{unique\_input\_addresses} и \textit{unique\_output\_addresses}~--- как практически линейно зависимые от агрегированных счётчиков входов и выходов,~--- а также два дополнительных признака, выявленных алгоритмом попарной проверки. Итоговое признаковое пространство составляет~30 переменных.

\subsubsection{Разбиение выборки}

Датасет разбивается на три непересекающихся подмножества: обучающую выборку~(70\%), валидационную~(15\%) и тестовую~(15\%). Разбиение выполняется стратифицированно~--- пропорции классов сохраняются во всех трёх частях. Это критически важно при дисбалансе: нестратифицированное разбиение может случайно недопредставить миноритарный класс~(\textit{mining\_pool},~89 образцов) в тестовой выборке, что сделает оценку качества нестабильной. После стратификации в тестовой выборке находится~$\lceil 89 \cdot 0{,}15 \rceil = 14$~адресов класса~\textit{mining\_pool}.

Важно, что оверсэмплинг и нормализация применяются исключительно к обучающей выборке, тогда как валидационная и тестовая выборки остаются в исходном виде~--- это исключает утечку данных (data leakage) и обеспечивает объективную оценку обобщающей способности модели.

\subsubsection{Устранение дисбаланса классов}

Дисбаланс классов с коэффициентом $\approx$~39,85 (отношение числа образцов мажоритарного класса к миноритарному) является существенной проблемой: классификатор, обученный на несбалансированных данных, стремится максимизировать суммарную точность за счёт занижения предсказаний для редких классов. Для устранения дисбаланса рассматриваются два метода оверсэмплинга:

\textbf{SMOTE}~(\acrlong{smote})~\cite{smote_paper}~--- для каждого образца миноритарного класса $\mathbf{x}_i$ выбирается $k$ ближайших соседей того же класса~(по умолчанию $k=5$). Синтетический образец $\tilde{\mathbf{x}}$ формируется интерполяцией:
\[
    \tilde{\mathbf{x}} = \mathbf{x}_i + \lambda \cdot (\mathbf{x}_{nn} - \mathbf{x}_i), \quad \lambda \sim U(0, 1),
\]
где $\tilde{\mathbf{x}}$~--- синтетический образец; $\mathbf{x}_i$~--- исходный образец миноритарного класса; $\mathbf{x}_{nn}$~--- случайно выбранный из $k$ ближайших соседей образец того же класса; $\lambda \sim U(0,1)$~--- коэффициент интерполяции. Такой подход создаёт разнообразные, но не выходящие за пределы выпуклой оболочки класса образцы.

\textbf{ADASYN}~(\acrlong{adasyn})~--- адаптивная модификация SMOTE, при которой интенсивность генерации для каждого образца пропорциональна плотности ошибок в его окрестности: больше синтетических примеров генерируется там, где класс труднее всего отделить от мажоритарных. Формально для образца $\mathbf{x}_i$ вычисляется $r_i = \Delta_i / K$, где $\Delta_i$~--- число соседей из мажоритарного класса среди $K$ ближайших соседей образца $\mathbf{x}_i$~--- доля таких соседей; число синтетических образцов для $\mathbf{x}_i$ пропорционально $\hat{r}_i = r_i / \sum_j r_j$, где суммирование ведётся по всем образцам миноритарного класса.

Выбор между SMOTE и ADASYN производится автоматически: прокси-модель Random Forest~(50 деревьев) обучается раздельно на каждом варианте оверсэмплинга и оценивается по взвешенной F1-мере на оригинальной~(не синтетической) валидационной выборке. По результатам сравнения выбран SMOTE~(val~F1~=~0,7475 против 0,7337 у ADASYN). Предпочтение SMOTE объясняется тем, что ADASYN, концентрируя генерацию в граничных зонах, увеличивает шум в разделяющей гиперплоскости~--- это может ухудшать обобщение. После применения SMOTE обучающая выборка расширяется до~14\,892~образцов с равным представлением каждого класса.

Дополнительно модели Random Forest и LightGBM обучаются с параметром \texttt{class\_weight="balanced"}, который масштабирует веса ошибок обратно пропорционально частоте класса. Это приводит к тому, что ошибка на одном образце класса \textit{mining\_pool} штрафуется примерно в 40 раз сильнее, чем ошибка на образце класса \textit{miner}.

\subsubsection{Стандартизация признаков}

После оверсэмплинга признаки стандартизируются методом \textit{StandardScaler}: каждый признак $x_j$ преобразуется к нулевому среднему и единичному стандартному отклонению:
\[
    \hat{x}_j = \frac{x_j - \mu_j}{\sigma_j},
\]
где $x_j$~--- исходное значение $j$-го признака; $\hat{x}_j$~--- его стандартизованное значение; $\mu_j$ и $\sigma_j$~--- среднее и стандартное отклонение признака $j$, вычисленные по обучающей выборке~\cite{scikit_learn}. Стандартизация необходима для MLP и kNN, поскольку оба метода чувствительны к масштабу признаков; для деревьевидных методов она не влияет на результат, однако сохраняется для унификации пайплайна.

\subsection{Обучение MLP}

Полносвязная нейронная сеть~(\acrshort{mlp}) рассматривается как основная нейросетевая модель в данной работе. MLP~--- это направленный граф слоёв нейронов: входной слой принимает вектор стандартизованных признаков размерности~30; далее следуют три скрытых слоя~(256, 128 и 64 нейрона); выходной слой содержит 6 нейронов~--- по одному на класс~--- с функцией активации~\textit{Softmax}, преобразующей логиты в вектор вероятностей.

На каждом скрытом слое применяется функция активации~\acrshort{relu}:
\[
    \text{ReLU}(z) = \max(0, z),
\]
где $z$~--- взвешенная сумма входов нейрона. Функция обеспечивает нелинейность преобразования и устраняет проблему затухающих градиентов. Обучение производится методом обратного распространения ошибки: вычисляется градиент функции потерь по всем весам сети, после чего веса обновляются оптимизатором Adam~\cite{adam_optimizer}~--- адаптивным стохастическим градиентным спуском с моментом первого и второго порядка:
\[
    \theta_{t+1} = \theta_t - \frac{\eta}{\sqrt{\hat{v}_t} + \varepsilon} \cdot \hat{m}_t,
\]
где $\theta_t$~--- вектор параметров модели на шаге $t$; $\hat{m}_t$ и $\hat{v}_t$~--- скорректированные на смещение оценки первого и второго момента градиента; $\eta$~--- скорость обучения; $\varepsilon$~--- малая константа для численной стабильности. В качестве функции потерь используется кросс-энтропия.

Для предотвращения переобучения применяется L2-регуляризация~(\textit{weight decay}) с коэффициентом~$\alpha = 10^{-4}$: к функции потерь добавляется штрафной член $\alpha \|\mathbf{W}\|_F^2$, препятствующий росту весов. Дополнительно используется ранняя остановка (\textit{early stopping}): если взвешенная F1-мера на валидационной доле обучающей выборки~(10\%) не улучшается на протяжении 10~итераций подряд, обучение прекращается и восстанавливаются веса с наилучшим значением метрики.

Гиперпараметры MLP зафиксированы по результатам предварительных экспериментов без использования автоматического поиска: объём обучающей выборки~(14\,892~образцов) в сочетании с высокой сложностью MLP делает перебор конфигураций через кросс-валидацию нецелесообразным по времени. Итоговые параметры: скорость обучения~$5 \cdot 10^{-4}$, размер батча~32, $\alpha = 10^{-4}$, максимальное число итераций~200.

\subsection{Ансамблевые методы}

\subsubsection{Random Forest}

Random Forest~\cite{scikit_learn}~--- ансамблевый метод, основанный на принципе бэггинга~(\textit{bootstrap aggregating}): для каждого из $B$ деревьев формируется бутстрэп-подвыборка~(с возвращением) из обучающей выборки, а при выборе оптимального разбиения в каждом узле рассматривается случайное подмножество признаков размером~$\sqrt{p}$ или~$\log_2 p$. Итоговый класс определяется мажоритарным голосованием деревьев. Случайность при отборе признаков снижает корреляцию между деревьями и повышает устойчивость ансамбля к переобучению.

Для борьбы с дисбалансом классов используется параметр \texttt{class\_weight="balanced"}, при котором вес каждого класса $c$ задаётся как $w_c = n / (K \cdot n_c)$, где $n$~--- полное число образцов, $K$~--- число классов, $n_c$~--- число образцов класса $c$.

Пространство поиска гиперпараметров:
\begin{itemize}
    \item число деревьев $B$: $\{100, 200\}$;
    \item максимальная глубина дерева: $\{\text{без ограничения},\; 15\}$;
    \item минимальное число образцов для разбиения узла: $\{2,\; 5\}$;
    \item число признаков при разбиении: $\{\sqrt{p},\; \log_2 p\}$.
\end{itemize}

\subsubsection{XGBoost}

XGBoost~\cite{xgboost_paper}~--- реализация градиентного бустинга на деревьях решений. В отличие от бэггинга, бустинг строит деревья последовательно: каждое новое дерево $h_t$ обучается аппроксимировать отрицательный градиент функции потерь $\mathcal{L}$ по предсказаниям текущего ансамбля. Итоговый ансамбль~--- аддитивная модель:
\[
    F_T(\mathbf{x}) = \sum_{t=1}^{T} \eta \cdot h_t(\mathbf{x}),
\]
где $F_T(\mathbf{x})$~--- итоговое предсказание ансамбля из $T$ деревьев для объекта $\mathbf{x}$; $h_t(\mathbf{x})$~--- предсказание $t$-го дерева; $\eta$~--- скорость обучения, масштабирующая вклад каждого дерева.

Ключевое отличие XGBoost~--- включение в целевую функцию регуляризационных членов, штрафующих сложность каждого дерева:
\[
    \widetilde{\mathcal{L}}^{(t)} = \sum_{i} l\bigl(y_i,\, \hat{y}_i^{(t)}\bigr) + \gamma T_t + \frac{1}{2}\lambda \|\mathbf{w}_t\|^2,
\]
где $\widetilde{\mathcal{L}}^{(t)}$~--- регуляризованная целевая функция на шаге $t$; $l(y_i, \hat{y}_i^{(t)})$~--- функция потерь для $i$-го обучающего образца; $y_i$~--- истинная метка класса; $\hat{y}_i^{(t)}$~--- текущее предсказание ансамбля; $T_t$~--- число листьев добавляемого дерева; $\mathbf{w}_t$~--- вектор весов листьев; $\gamma$ и $\lambda$~--- гиперпараметры регуляризации. Это предотвращает переобучение даже на глубоких деревьях.

Пространство поиска гиперпараметров:
\begin{itemize}
    \item число деревьев $T$: $\{100, 200\}$;
    \item максимальная глубина дерева: $\{4,\; 6\}$;
    \item скорость обучения $\eta$: $\{0{,}1;\; 0{,}2\}$;
    \item доля случайно выбираемых образцов для дерева~(\textit{subsample}): $\{0{,}8;\; 1{,}0\}$;
    \item доля случайно выбираемых признаков~(\textit{colsample\_bytree}): $\{0{,}8;\; 1{,}0\}$.
\end{itemize}

Параметры \textit{subsample} и \textit{colsample\_bytree} вносят стохастичность, аналогичную бэггингу, и снижают корреляцию между деревьями.

\subsubsection{LightGBM}

LightGBM~\cite{lightgbm_paper}~--- высокопроизводительная реализация градиентного бустинга, использующая два ключевых нововведения. Во-первых, \textit{алгоритм листового роста}: вместо уровневого (depth-wise) наращивания дерева LightGBM на каждом шаге добавляет лист к узлу с наибольшим приростом функции. Это позволяет строить глубокие несимметричные деревья за меньшее число шагов. Во-вторых, \textit{гистограммная аппроксимация}: непрерывные признаки дискретизируются в $B$ бинов~(по умолчанию 255), что снижает вычислительную сложность поиска оптимального разбиения с $O(n \cdot p)$ до $O(B \cdot p)$. В совокупности это обеспечивает скорость обучения в 10--100 раз выше, чем у XGBoost, при сопоставимом качестве.

Используется параметр \texttt{class\_weight="balanced"}. Пространство поиска:
\begin{itemize}
    \item число деревьев: $\{100, 200\}$;
    \item максимальная глубина дерева: $\{4,\; 6,\; -1\}$~(без ограничения);
    \item скорость обучения: $\{0{,}1;\; 0{,}2\}$;
    \item максимальное число листьев~(\textit{num\_leaves}): $\{31,\; 63\}$;
    \item доля образцов для построения дерева~(\textit{subsample}): $\{0{,}8;\; 1{,}0\}$.
\end{itemize}

Параметр \textit{num\_leaves} является основным регулятором сложности в LightGBM: при листовом росте глубина дерева сама по себе не ограничивает сложность, поэтому ограничение числа листьев эквивалентно ограничению глубины в XGBoost.

\subsubsection{Stacking}

Stacking~(\textit{StackingClassifier})~--- двухуровневый ансамбль. На первом уровне три базовых модели~--- Random Forest, XGBoost и LightGBM~--- обучаются на обучающей выборке. Для формирования признаков мета-уровня применяется трёхкратная стратифицированная кросс-валидация: обучающая выборка делится на три фолда; каждая базовая модель обучается на двух фолдах и предсказывает вероятности для третьего. Таким образом для каждого объекта обучающей выборки вычисляется вектор вероятностей без использования самого объекта в обучении~--- это исключает утечку данных. Итоговая матрица мета-признаков имеет размерность $n_{\text{train}} \times (3 \cdot K)$, где~$K = 6$~--- число классов.

На втором уровне мета-классификатор~--- логистическая регрессия~--- обучается на мета-признаках. При оценке на тестовой выборке базовые модели дообучаются на полной обучающей выборке и формируют мета-признаки для тестовых объектов; затем мета-классификатор выдаёт финальное предсказание.

\subsubsection{kNN}

kNN~--- непараметрический классификатор, не требующий явного обучения. Для классификации нового объекта $\mathbf{x}_{\text{test}}$ вычисляется евклидово расстояние до всех объектов обучающей выборки, отбираются $k$ ближайших, и класс определяется большинством голосов:
\[
    \hat{y} = \arg\max_{c} \sum_{i \in N_k(\mathbf{x}_{\text{test}})} \mathbf{1}[y_i = c],
\]
где $\hat{y}$~--- предсказанный класс; $c$~--- перебираемый класс; $N_k(\mathbf{x}_{\text{test}})$~--- множество индексов $k$ ближайших соседей объекта $\mathbf{x}_{\text{test}}$; $y_i$~--- истинная метка $i$-го соседа; $\mathbf{1}[\cdot]$~--- индикаторная функция. Используется $k=7$ без дополнительной настройки~--- как простейший базовый классификатор. Метод чувствителен к масштабу признаков~(поэтому необходима стандартизация) и к «проклятию размерности» в высокодимензиональных пространствах, однако хорошо иллюстрирует, насколько линейно не разделимы классы в исходном признаковом пространстве.

\subsection{Настройка гиперпараметров}

Для Random Forest, XGBoost и LightGBM применяется случайный поиск гиперпараметров (\textit{RandomizedSearchCV})~\cite{scikit_learn}. В отличие от полного перебора (\textit{GridSearchCV}), случайный поиск равномерно сэмплирует заданное число ($n_{\text{iter}}$) конфигураций из декартова произведения пространств параметров, что позволяет существенно сократить время поиска при сохранении вероятности найти хорошее решение~--- особенно если оптимум сосредоточен в небольшой подобласти.

Оценка каждой конфигурации производится методом стратифицированной кросс-валидации~(\textit{StratifiedKFold}, $n=2$): обучающая выборка делится на два стратифицированных фолда, модель обучается на одном и оценивается на другом~--- поочерёдно; результирующий балл~--- среднее значение взвешенной F1-меры по двум фолдам. Выбирается конфигурация с наибольшим средним баллом.

Число итераций поиска: для Random Forest~--- 4, для XGBoost и LightGBM~--- 6. Внешний параллелизм~(\texttt{n\_jobs=-1} в~\textit{RandomizedSearchCV}) позволяет проверять все конфигурации параллельно; внутри каждой модели используется \texttt{n\_jobs=1} для предотвращения конкурентного захвата потоков при вложенном параллелизме.
